%! AP Statistics — Unit 9, Lesson 9.1 — Homework (Answer Key)
\documentclass[10pt]{article}
\usepackage{apstats-article}
\usepackage{apstats-key}

\begin{document}

\pageheader{Unit 9, Lesson 9.1}{Homework --- Answer Key}
\namedateperiod

\vspace{0.1in}

\textbf{Problem 1.} A student fits a least-squares regression line to data on daily
temperature ($^\circ$F, $x$) and the number of hot drinks sold at a caf\'e ($y$) from a
random sample of 20 days. The equation is $\hat{y} = 85 - 1.4x$.

\begin{enumerate}[label=\alph*.]

  \item Interpret the slope in context.

        \ans{For each 1$^\circ$F increase in daily temperature, the predicted number of
        hot drinks sold decreases by 1.4.}

  \item The student concludes: ``Temperature causes fewer hot drinks to be sold because
        the slope is negative.'' Identify \textbf{two} problems with this conclusion.

        Problem 1: \ans{Correlation does not imply causation. The negative slope shows
        an association between temperature and hot drink sales, but does not establish
        that temperature \textit{causes} the change in sales. A confounding variable
        (e.g., season) could explain both.}

        Problem 2: \ans{The sample slope $b = -1.4$ may not reflect the true population
        slope $\beta$. If $\beta = 0$ (no real relationship), random variation in a
        sample of 20 days could still produce $b = -1.4$. We cannot conclude $\beta \ne 0$
        without a significance test.}

  \item A second study using $n = 200$ days finds $b = -1.4$ as well. In which study
        would the slope be more convincing evidence of a real relationship? Explain.

        \ans{The study with $n = 200$ would be more convincing. Larger sample sizes
        reduce sampling variability in $b$, so obtaining $b = -1.4$ from 200 days is
        much less likely to be due to chance alone than from just 20 days. The same
        slope value is stronger evidence against $\beta = 0$ when $n$ is large.}

\end{enumerate}

\vspace{0.15in}

\textbf{Problem 2.} For each scatter plot description below, state the question a
statistician would ask about the population.

\begin{enumerate}[label=\alph*.]

  \item A scatter plot of shoe size ($x$) vs.\ height (inches, $y$) for 50 adults shows
        a positive trend with $b = 0.62$.

        \ans{$H_0\colon \beta = 0$ --- there is no linear relationship between shoe size
        and height for adults in the population. The statistician asks: is the observed
        slope $b = 0.62$ convincingly different from 0, or could it have arisen by
        chance even if $\beta = 0$?}

  \item A scatter plot of study hours ($x$) vs.\ exam score ($y$) for 8 students shows
        $b = 4.5$, but the points are very scattered around the line.

        \ans{$H_0\colon \beta = 0$ --- there is no linear relationship between study
        hours and exam score in the population. With only $n = 8$ and high scatter, the
        slope $b = 4.5$ may not be statistically significant; a test is needed to
        determine whether this $b$ is larger than expected by chance.}

  \item A scatter plot of screen time (hours, $x$) vs.\ sleep duration (hours, $y$) for
        100 teens shows $b = -0.15$ with little scatter around the line.

        \ans{$H_0\colon \beta = 0$ --- there is no linear relationship between screen
        time and sleep duration in the population. With $n = 100$ and little scatter,
        the slope $b = -0.15$ is more likely to represent a real effect, but a
        significance test is still needed to formally conclude $\beta \ne 0$.}

\end{enumerate}

\end{document}
